Graph neural networks (GNNs) are usually motivated by homophily: neighbouring nodes tend to share labels, so averaging over neighbours denoises features. The graph convolutional network (GCN) of Kipf and Welling~\cite{kipf2016semi} is the canonical example, and a line of work on heterophily argues that it should fail when neighbours have different labels and proposes designs such as separating the ego embedding from the neighbour embedding~\cite{zhu2020beyond} or learning signed propagation weights~\cite{chien2020adaptive}. Yet the evidence is mixed. Ma et al.\ show that GCN can do well on some heterophilous graphs~\cite{ma2021is}, Luan et al.\ argue that homophily is the wrong summary statistic and that post-aggregation node similarity matters more~\cite{luan2022revisiting}, and Platonov et al.\ find that standard GNNs usually beat specialised heterophily models once benchmark flaws are fixed~\cite{platonov2023critical}.

On real benchmarks, homophily is confounded with feature quality, degree, graph size and split protocol, which makes it hard to say which factor drives a result. A contextual stochastic block model (CSBM)~\cite{deshpande2018contextual} removes the confound: it generates a graph with a planted community structure together with Gaussian node covariates, so one can vary edge homophily and feature informativeness independently. Theory for this family shows that graph convolution widens the regime in which the classes are linearly separable under homophily~\cite{baranwal2021graph}; it is less clear what happens in between, where the graph is nearly uninformative and averaging mixes classes.

This paper is a small, controlled empirical map of that question. We sweep homophily from 0 to 1 and feature strength over three levels, and compare an MLP, a 2-layer GCN, a GCN with separate self and neighbour weights, and label propagation, all written from scratch in torch. We register five hypotheses before running the main grid (Section~\ref{sec:setup}) and report each verdict, including the ones that fail. Our contributions are: (i) a win-region map over 33 homophily$\times$informativeness cells with 5 seeds each; (ii) a test of the claim that GCN drops below the MLP at mid homophily, which is supported only partially and localised below the middle; (iii) an ablation showing that the ego/neighbour separation, not extra capacity, is what matters under heterophily; and (iv) a mean-degree sensitivity check at one cell. The study is deliberately small (900 nodes, 2 layers, CPU, 5 seeds), and we state what it cannot show in Section~\ref{sec:limitations}.
